\subsection{The exact coefficient period left by a general form-one repair}

\paragraph{A-S09-16 (Proved; maximal simultaneous repair and its remaining period).}
Fix a positive weight and identify its coefficient group with
$\mathbb Z$.  Let $X$ be a simplicial nonnegative cochain complex
with $X_0=\mathbb Z[0]$ and $H^0(X_n)=\mathbb Z$ through the
coefficient map in every simplex.  Let $f:X\to A$, $g:X\to B$
preserve coefficients, where both targets have $H^0=\mathbb Z$
and $H^1=0$ in every simplex.  Put
\[
 \Phi_n=\mathbb Z^{\{0,\ldots,n\}},\qquad
 L=\Phi/\underline{\mathbb Z},
 \tag{S09.53}
\]
with simplicial operators given by restriction of vertex
functions and with the denominator consisting of constant
functions.  There is a canonical simplicial homomorphism
\[
 \Delta:H^1(X)\longrightarrow L                       \tag{S09.54}
\]
and an explicit simplicial cochain extension $X\hookrightarrow E$
with unchanged terms in positive cochain degrees, extending both
target maps, such that
\[
 H^0(E)=\underline{\mathbb Z},\qquad
 H^1(E)=\operatorname{im}\Delta\subset L,
 \qquad H^q(E)=H^q(X)\quad(q\geq2).                    \tag{S09.55}
\]
The degree-zero groups of $E$ are free whenever $X^1$ is
termwise free.  In every coefficient-preserving extension
of the prescribed pair of maps into a complex with
$H^1=0$ and simplex-zero $H^0=\mathbb Z$, the homomorphism
(S09.54) vanishes.  Thus (S09.54) is the exact obstruction
for simultaneous degree-zero repair; no remaining
linear obstruction is suppressed.

\begin{proof}
For $c\in Z^1(X_n)$ choose a primitive $a\in A_n^0$ of
$f(c)$.  Every vertex of $c$ is zero, so every vertex of
$a$ is a coefficient cycle.  Define
\[
 \delta_f(c)=[(a_0,\ldots,a_n)]\in L_n.               \tag{S09.56}
\]
The difference between any two primitives is a constant
coefficient cycle in $A_n^0$.  Consequently (S09.56) is
independent of the choice, is additive, and commutes
with every simplicial operator.  The same construction
gives $\delta_g$.  For $b\in X_n^0$ we may use $f(b)$
and $g(b)$ themselves as primitives.  Their vertex
values are the same coefficients $b_i$, so
\[
        \delta_f(db)=\delta_g(db)=[(b_0,\ldots,b_n)].
 \tag{S09.57}
\]
The difference $\delta_f-\delta_g$ therefore factors
through vertical cohomology, giving (S09.54).

Let $C\subset Z^1(X)$ be the kernel of this difference
before passing to cohomology.  On $C$, the two period
maps agree.  Define a pullback in simplicial abelian
groups by
\[
 E^0=C\mathop{\times}_{L}\Phi,
 \qquad
 E_n^0=\{(c,v):c\in C_n,\ v\in\Phi_n,
                          [v]=\delta_f(c)\},
 \qquad d(c,v)=c.                                    \tag{S09.58}
\]
In positive cochain degrees retain $X$ and its differential.
This is a complex because every $c$ is closed.  The
inclusion in degree zero is
\[
                 b\longmapsto (db,(b_0,\ldots,b_n)). \tag{S09.59}
\]
It is injective: an element in its kernel is a closed
zero-form with all vertex values zero, hence is the zero
coefficient.  The construction is simplicial without
making a vertex choice.

For $(c,v)\in E_n^0$, there is exactly one primitive
$a\in A_n^0$ of $f(c)$ whose vector of vertex values
is $v$.  Existence follows from the equality of their
classes in $L_n$: adjust any primitive by a constant
coefficient.  Uniqueness follows from $H^0(A_n)=\mathbb Z$.
Sending $(c,v)$ to this primitive extends $f$ and is
additive and simplicial by uniqueness.  Since
$\delta_f(c)=\delta_g(c)$, the same construction
extends $g$.  This argument uses no flatness of the
targets.

The kernel of $d:E_n^0\to X_n^1$ consists exactly of
the pairs $(0,v)$ with $v$ constant, proving the
$H^0$ statement.  The map $\Phi_n\to L_n$ is
surjective, so $dE_n^0=C_n$.  Hence
$H^1(E_n)=Z^1(X_n)/C_n=\operatorname{im}\Delta_n$.
All higher cohomology groups are unchanged.
Finally $E_n^0$ is a subgroup of $X_n^1\oplus\Phi_n$,
which proves freeness under the stated hypothesis.

To prove the asserted necessity, let $j:X\to Y$
extend the two prescribed maps, and let $c\in Z^1(X_n)$
become a boundary $dh=j(c)$ in $Y_n$.  Every vertex
of $h$ is a coefficient cycle because every vertex of
$c$ is zero.  The two target images of $h$ consequently
have the same vector of vertex coefficients.  They
are primitives of $f(c)$ and $g(c)$, so
$\delta_f(c)=\delta_g(c)$.  If $H^1(Y_n)=0$, this
holds for every $c$.  This proves necessity even
when $j$ is not injective and when new positive-degree
terms were added to $Y$.
\end{proof}

The obstruction can also be read on edges.  Restricting
a vertex vector to the edge from $0$ to $i$ recovers
$a_i-a_0$, so the whole map (S09.54) is determined by
the difference of the two endpoint periods in simplex
degree one.  For an admissible target in a sufficiently
large fixed weight, $\delta_f$ is the connecting map
induced by the exact sequence
$0\to\underline{\mathbb Z}\to A^0\to Z^1(A)\to0$,
expressed by actual primitives and vertex restrictions.
This gives a direct coefficient transgression, including
its signs, rather than an unspecified secondary operation.

\paragraph{The precise PD attachment at this one weight.}
Let $D$ be a strict PD cochain algebra with weight-zero
part $\mathbb Z$, and apply (S09.58) to $X=D^w$.
The degree-zero module $D^{w,0}$ and its extension
$E^0$ lie in the positive PD ideal.  The pushout of
graded PD algebras
\[
 D'=D\mathop{\otimes}_{\Gamma(D^{w,0})}\Gamma(E^0)
 \tag{S09.60}
\]
has the special differential extending the differential
of $D$ and sending $(c,v)$ to $c$.  Formula (S09.59)
ensures that this differential respects the amalgamation
relations.  On a divided power it is the required
$d\gamma_m(e)=\gamma_{m-1}(e)\,de$; its square is
zero since $dc=0$ and odd squares vanish.  The two
target maps constructed above extend over (S09.60)
by the PD universal property.  At weight $w$ the
complex is exactly $E$: all new products of positive
weight and all higher divided powers occur in larger
weights.  This proves the current-weight PD step.

There is a precise freeness statement, without
assuming that the pushout has no later relations.
In each simplex,
\[
 E_n^0/D_n^{w,0}
       \cong\ker\bigl(\Delta_n:H^1(D_n^w)\to L_n\bigr).
 \tag{S09.61}
\]
Indeed the projection $(c,v)\mapsto[c]$ has this
image, and its kernel is (S09.59), using the coefficient
adjustment when $c=db$.  When the group in (S09.61)
is free, choose a splitting in that simplex.
The underlying PD algebra (S09.60) is then
$D_n\otimes\Gamma(E_n^0/D_n^{w,0})$, so it is
termwise free whenever $D_n$ is.  A simplicial splitting
is unnecessary: the pushout and its target maps were
defined before these choices.  The assertion in
(S09.61) applies, in particular, when $H^1(D_n^w)$
is free.  Without this quotient hypothesis no freeness
of later weights is claimed.  In all cases the
construction alone makes no assertion about later
vertical or horizontal admissibility.

\paragraph{An actual nonzero period with $D_0=S$.}
Let $G$ be the standard weight-one PD de Rham theory,
with simplex coordinates $t_0,\ldots,t_n$ satisfying
$\sum t_i=T$.  Use $L$ in (S09.53), placed in form
degree one and weight one, and form
\[
 D=G\otimes_{\mathbb Z}\Lambda(L),\qquad dL=0.
 \tag{S09.62}
\]
The exterior algebra carries its canonical graded PD
structure; equivalently it is the free strict graded
PD algebra on this odd module.  Its simplicial groups
are free.  The following are actual strict PD cochain
maps, both equal to the identity on $G$:
\[
 f_0:D\to G,\quad [a_0,\ldots,a_n]\longmapsto0,
 \qquad
 f_1:D\to G,\quad [a_0,\ldots,a_n]
                              \longmapsto\sum_i a_i\,dt_i.
 \tag{S09.63}
\]
The second is well defined because $\sum_i dt_i=0$,
and it commutes with faces and degeneracies by their
coordinate formulas.  Its images are closed odd forms,
so it extends to the asserted strict PD map.

Here $D_0=S$, $H^0(D_n)=S$, and
\[
 H^1(D_n)^w= S^{w-1}\otimes L_n\quad(w\geq1).
 \tag{S09.64}
\]
This follows by tensoring the exact vertical complex
of $G_n$ with the free exterior algebra with zero
differential.  All groups are free and $q\leq w$.
Moreover every row $D^{w,q}$ is simplicially acyclic
for $w>q$: it is the finite direct sum of
$G^{w-m,q-m}\otimes\Lambda^m(L)$, $0\leq m\leq q$,
and $w-m>q-m$.  Each indicated row of $G$ is an
acyclic simplicial free abelian group.  The integral
Eilenberg--Zilber equivalence and its contracting
normalized complex make each displayed tensor
acyclic.

For $c=T^{[w-1]}[a_0,\ldots,a_n]$, a primitive of
its $f_1$-image is $T^{[w-1]}\sum_i a_i t_i$.
Its vertex values are $wa_iT^{[w]}$, whereas the
$f_0$-image has the zero primitive.  Consequently
the difference period on (S09.64), with the order
$\delta_{f_1}-\delta_{f_0}$, is exactly
\[
                       L\xrightarrow{\ w\ }L.        \tag{S09.65}
\]
It is nonzero in every positive weight.  In simplex
degree one, the generator $[0,1]$ has period $w$.
The necessity part of A-S09-16 therefore proves that
the prescribed pair (S09.63) has no extension through
any eventual admissible theory.  This remains true
for noninjective maps from $D$ into a proposed
extension.  The algebra $D$ is a partial source with
the explicitly displayed nonzero vertical $H^1$;
it is not itself an admissible object and gives no
counterexample to the main common-source assertion.
